\section{经济与社会：维持一张封建网络要付出什么}

西周没有留下可供我们逐户相加的税册，也不能简单想象成后世那种由统一市场或统一税制驱动的国家。现存材料所能看见的，是农业、田地、役使、赏赐、军旅、祭祀与宗族关系彼此扣连。王室要使远近封国继续承认其册命和裁断，就要有粮食、道路、工匠、可调动的人力，以及能把这些资源转成礼器和军功的组织。

这幅图像有明显的材料边界。金文常由受命者为纪念王命、功劳和家族祭祀而铸；它最清楚地照见贵族如何记录资源与身份，并不替无名的耕作者说话。周原、丰镐、成周的居址、仓储、作坊和墓地，是考古所见的物质条件，却不能自行给出赋役比例、人口总数或行政机构名称。至于\sourcebook{尚书}、\sourcebook{诗经}和\sourcebook{史记}，分别保存训诰、歌诗与后出的连贯叙事；它们能提出周人及后来人怎样理解秩序，不能补成一份西周现场的经济统计表。

\subsection{田地不是一张可以直接摊平的地图}

后世常用整齐的“井田制”来想象西周土地。现有的器铭却呈现另一种更难一笔带过的世界：田地、邑落、宅地、贝、马、丝帛以及役使关系，会在贵族之间被确认、给付或争议。它们说明资源权利能够被命名、见证和传给后代；但这些零散的精英记录不足以让我们复原全境土地分配，更不能据此断言人人都服从同一种田制。

\begin{event}{西周中期至晚期（具体王年有争议）}{裘卫诸器所见的资源给付与关系交换}{器铭可靠；交易性质和术语解释有争议}{周王畿贵族社会}\eventanchor{WZ-QIU-WEI-TRANSACTIONS}{西周·裘卫诸器的资源交换}
裘卫盉、簋、盘把贝、马、丝帛及与服务相连的资源写入器铭。它们让人看见，贵族关系并不只靠王室单向赏赐维持，也有需要被交割和追认的财物往来。把这种记录直接叫作现代市场买卖同样过头：器文保存的是一端精英关系，债务、役使、赠与和权利转移究竟怎样区分，仍有不同解释。

不过，正因为要把若干物项并列并留下凭据，资源流动已不能被简化为“贡纳”二字。五祀卫鼎、九年卫鼎和卫盉涉及田地、邑落与役属的交割；曶鼎以马、丝帛和田地处理给付；大克鼎与晚期克鼎又把田宅、贝和王命同家族祭祀连在一起。这些铭文所能确定的是权利与服务在贵族网络中相互嵌合，不能把器铭中的一个术语不加说明地换成近代的私有产权或雇佣劳动。
\end{event}

土地还会引起争执。五年、六年琱生簋所保留的确认，表明一次处置未必消除了后续解释空间；散氏盘则把田界、赔偿、见证和誓约写成可由后人追认的安排。这里可见的不是一部完备成文法，而是家族和相邻资源关系需要借仪式、见证与铭文固定。边界之所以值得郑重保存，正因为它关乎产出、役使与家族声望。

\begin{event}{西周晚期（约前九世纪；具体王年未定）}{散氏盘保存田界争议的处置}{器铭可靠；地望、程序与社会范围有争议}{周王畿贵族社会}\eventanchor{WZ-SANSHI-BOUNDARY-CASE}{西周·散氏盘的田界争议}
散氏盘把赔偿、田界和誓约并置，显示纠纷可被转换为一段郑重保存的文本。它不能证明西周每一处田地都按同一程序诉讼，却提醒我们：王命之外，贵族秩序也要靠协商、见证和可以跨代调用的凭据维持。把它同裘卫、卫、曶、琱生诸器合看，较稳妥的结论是资源权利并非静止不动；把它们叫作自由土地市场，则超出了材料。
\end{event}

\subsection{贡纳与劳役：供给不是抽象的背景}

王室的命令能抵达多远，取决于道路、仓储和人力能否把物资送到需要的地方。周原、丰镐的遗址显示聚落、窖藏、仓储和手工业遗存的集中；植物考古与铸铜遗存又提示，有些物资须被储备和集中处理。这是考古能给出的趋势，不是能够折算为年度粮食总量或赋役率的账簿。

\begin{event}{西周中晚期（为持续过程）}{王畿的聚落、仓储和手工业支撑王室动员}{考古趋势可靠；行政规模、赋役比例与人口数字不可确证}{宗周王畿及周原、丰镐地区}\eventanchor{WZ-WESTERN-CAPITAL-ECONOMY}{西周·王畿的仓储与手工业}
双都、祭祀、赏赐与征伐不可能凭一纸命令运行。周原的仓储、窖藏和作坊遗存，丰镐的居址与交通条件，连同金文和歌诗所见的王室活动，共同说明王畿需要稳定调集粮食、工匠和运输能力。它们没有标明由哪一户负担多少，也不许可我们把每一处遗存径直归作官仓；但足以使“劳役”和“供给”不再只是制度名称。

成周的道路、居住区和仓储遗存，则把\eventref{WZ-EAST-CHENGZHOU}{西周·营建成周}之后的东方经营落到物质节点上。王室要在东方举行朝会、转运物资并维系受命关系，便须有能够容纳人员和货物的城邑。传世\sourcebook{尚书·召诰}、\sourcebook{尚书·洛诰}和后出的\sourcebook{史记}叙事，提供了营建与安置的政治框架；何尊及遗址把某些王命和空间条件拉近当时。三类材料可以互证，却不能让我们写出迁置户数或每项劳役的定额。
\end{event}

贡纳也应作这样的辨析。宣王时期说的兮甲盘，确实把戍守、马匹与贡纳管理置于同一边务关系中；它可说明边防依赖日常登记、交运和责任分配，不能改写成一条覆盖整个西周的税法。军旅诗如《采薇》保存的是戍役和远征的经验性语言，成诗、编定与后世解释的层次却使它不能单独充当征发清单。所谓“地方承担成本”，在材料中首先表现为这些人、马、物资和道路必须被持续调动，而非可以精确核算的一笔总额。

\subsection{青铜：把资源变成可见的关系}

\begin{event}{西周中期（约前十世纪；为持续过程）}{青铜生产、原料调度与工匠组织扩大并区域化}{考古趋势可靠；产量、控制方式与具体年份不定}{周及各地贵族网络}\eventanchor{WZ-MID-BRONZE-PRODUCTION}{西周·青铜生产与贵族网络}
一件礼器不是孤立的奢侈品。铜、锡、燃料、范铸技术、工匠与运输要在相当长的时间内被组织起来，周原、成周及诸侯墓地所见器物和铸铜遗址，使这种供应链的存在可以被追问。器物铭文又把册命、征伐、赏赐和祭祀保存下来：资源经过工坊，成为一件可以在朝会、祭祀和家族传承中反复使用的政治记忆。

因此，青铜器既能说明集中能力，也能显示地方贵族的能力。王室需要礼器来使受命和功劳可见；受命家族也能借器物让其服务、祖先与资源基础被后代记住。考古所见的区域差别和不同墓地组合提醒我们，不宜把这套礼器语言想成由中心完全整齐复制的统一制度。金文最接近的仍是铸器者及其交往圈，不会直接告诉我们矿料从何地以何种条件取得，或工匠和农户怎样评价这种动员。
\end{event}

墓葬、车马和礼器组合使等级有了可见形状。周原、丰镐、成周及诸侯墓地的差异，连同铭文中的家族记忆，说明贵族身份、继承与祭祀被物质化；它们也保留了地域差异。后世把“周礼”叙述为整齐的典范，是另一层历史记忆，不能倒过来抹平考古所见的不均衡。普通人的劳动构成这套秩序的条件，却很少以自己的名字进入这些器物和墓葬。

\subsection{王都离开关中以后}

\eventref{WZ-0771-HAOJING}{西周·镐京陷落}并不只是一座都城的军事失守。晚期周原层位和扶风地区调查可见部分聚落、窖藏与生产活动的中断或转变；关中聚落在东迁后也有调整。不过，遗址分期不能自动把每一处变化归因于前771年，更不能逐村追踪人群的去向。传世\sourcebook{史记·周本纪}叙述镐京危机和东迁，考古材料显示西部基础的变化节奏；二者相接处，正是不能用单一灾变故事填满的部分。

\begin{turningpoint}[制度得失：分封网络的效率与代价]
\term{它解决了什么}：在交通和文书工具有限的条件下，王室无法以少量官员直接经营广阔地区。封国与受命家族以本地宗族、土地、劳力和礼器网络承担守卫、经营与联络；王畿的仓储、作坊和道路则让王室仍能组织朝会、赏赐与征伐。

\term{谁承担代价}：器铭让我们看见的是得到田宅、贝、马和名分的贵族，遗址让我们看见的是仓储、作坊和聚落的物质条件。耕作、运输、筑城、戍守和从军者多半没有留下可辨姓名。可以确认的是动员离不开他们，不能确认的则是他们各自承担了多少赋役、怎样感受这一秩序。

\term{它留下什么压力}：以土地、役使和物资确认服务，能使王室命令落到地方，也让受命家族把资源、礼器和声望积累为代际基础。资源权利需要不断经由王命、见证和协商重申，边防又不断消耗人马和转运能力。到镐京危机与东迁时，关中的聚落和供给体系已开始以不同速度调整；原先连接中心与地方的渠道，因而也可能成为地方各自调动资源的渠道。
\end{turningpoint}
